% Generated draft: requires mathematical and source review.
\documentclass[11pt]{article}
\usepackage[margin=25mm]{geometry}
\usepackage{fontspec}
\setmainfont{Latin Modern Roman}
\usepackage{amsmath,amssymb,mathtools,unicode-math}
\setmathfont{Latin Modern Math}
\usepackage{enumitem,needspace,xurl,hyperref}
\setlength{\parindent}{0pt}
\setlength{\parskip}{5pt}
\setlength{\emergencystretch}{3em}
\newcommand{\N}{\mathbb{N}}
\newcommand{\Z}{\mathbb{Z}}
\newcommand{\Q}{\mathbb{Q}}
\newcommand{\R}{\mathbb{R}}
\newcommand{\C}{\mathbb{C}}
\newcommand{\F}{\mathbb{F}}
\newcommand{\A}{\mathbb{A}}
\newcommand{\PP}{\mathbb{P}}
\newcommand{\E}{\mathbb{E}}
\newcommand{\eps}{\varepsilon}
\newcommand{\dd}{\,\mathrm{d}}
\begin{document}
\section*{UnsolvedMath ID: 701526}
\textbf{Source problem code:} catalogue entry
\subsection{Definitions and mathematical statement}
A group is periodic if every element has finite order. A group is locally finite if each subgroup generated by finitely many elements is finite. An involution is an element of order exactly \(2\). A \(2\)-group here means a group in which every element has order a power of \(2\), allowing the identity to have order \(2^0=1\).

Let \(G\leq\operatorname{Sym}(\Omega)\) be periodic and two-transitive: for any \(x_1\ne x_2\) and \(y_1\ne y_2\) in \(\Omega\), there exists \(g\in G\) such that \(g(x_i)=y_i\) for \(i=1,2\). Suppose also that every three pairwise distinct points have trivial pointwise stabilizer:
\[
\{g\in G:g(x)=x,\ g(y)=y,\ g(z)=z\}=\{1\}
\qquad(x,y,z\in\Omega\text{ pairwise distinct}).
\]
These are the Zassenhaus-group conditions specified in the supplied record.

For a point \(\alpha\in\Omega\), let
\[
H=G_\alpha=\{g\in G:g(\alpha)=\alpha\}.
\]
Assume that \(H\) is a Frobenius group with locally finite Frobenius kernel \(U\trianglelefteq H\), and that \(U\) contains an involution. Here a Frobenius group is understood in the usual sense of a faithful transitive permutation group with nontrivial point stabilizers and trivial stabilizers of any two distinct points. Its Frobenius kernel consists of the identity together with the elements fixing no point in that Frobenius action; the hypothesis includes that this is the subgroup \(U\). This Frobenius action of \(H\) is distinct from the given action of \(G\) on \(\Omega\).

Must \(U\) be a \(2\)-group? Explicitly, must
\[
\forall u\in U\ \exists n\in\Z_{\geq 0}\quad \operatorname{ord}(u)=2^n
\]
hold? Since \(U\) is locally finite, this is equivalently the assertion that every finitely generated subgroup of \(U\) is a finite group of order a power of \(2\).

The supplied record attributes the finite-group case to Feit, but supplies no bibliographic details for that result.
\subsection{Short English statement}
N. M. Suchkov asks whether the locally finite Frobenius kernel of a point stabilizer in a periodic Zassenhaus group must be a \(2\)-group whenever it contains an involution. The record attributes the finite-group case to Feit and labels the general problem open; these are supplied source assertions, not independently verified assessments.
\subsection{Sources}
\begin{itemize}
\item Local UnsolvedMath registry, permanent ID \(701526\), exact source problem code \texttt{catalogue entry}, title \emph{Kourovka Notebook Problem 18.97}. The supplied record attributes the problem to N. M. Suchkov, proposed in 2014.
\item \emph{The Kourovka Notebook (Unsolved Problems in Group Theory)}, Problem 18.97, 18th issue (2014); the supplied background also identifies the 21st issue edition (2026), \texttt{arXiv:1401.0300v48}. The supplied collection description names E. I. Khukhro and V. D. Mazurov as editors.
\item The supplied statement credits Feit with proving the assertion for finite groups. No separate paper, theorem reference, or authoritative citation for that attribution is supplied.
\item The local record labels the problem \texttt{open} and records an update timestamp of \texttt{2026-10-16T14:00:00Z}. This registry status is supplied curation and has not been independently verified.
\end{itemize}
\end{document}
